% GENERATED FILE. Edit canonical lessons, then run npm run generate:books.
% canonical-source-hash: fnv1a64:979c8920145df4c9
% canonical-lessons: IT-C507-giardinaggio, IT-C507-seminare, IT-C507-annaffiare, IT-C507-rastrello, IT-C507-zappare

\chapter{Il giardinaggio}
\label{ch:it-il-giardinaggio}
\hlchaptermodality{507}

\begin{chapteropening}
\textbf{By the end of this chapter:} \emph{I can talk about gardening.}

Complete the last lesson of chapter 507: I can talk about gardening.
\end{chapteropening}

\section[il giardinaggio]{il giardinaggio — gardening}

\label{lesson:IT-C507-giardinaggio}



\begin{quote}
Before the new one: say the Italian for roll of film, then the Italian for to develop.
\end{quote}

\subsection*{You'll want to know: il giardinaggio}
\textbf{il giardinaggio} — \textquotedblleft{}gardening\textquotedblright{}.

\begin{quote}
Mia madre fa giardinaggio tutti i fine settimana.
\end{quote}

\begin{quote}
\emph{My mother does gardening every weekend.}
\end{quote}

\emph{Fare giardinaggio} — to garden.

\begin{cousinweb}
\emph{giardinaggio} comes from French \emph{jardinage}, built on \emph{jardin} \textquotedblleft{}garden\textquotedblright{}, from Medieval Latin \emph{jardīnus} (EtymDB), which is usually traced to a Frankish word for an enclosure.
\end{cousinweb}

\begin{grammarlens}[title={What you've built}]
One more word for telling, explaining and describing at B1.
\end{grammarlens}

\subsection*{Your turn}
\begin{itemize}
\raggedright
  \item \emph{Say it:} \emph{il giardinaggio}
  \item \emph{Say it:} \emph{il giardinaggio}, once more
  \item \emph{Say it:} a sentence of your own with \emph{il giardinaggio}
  \item \emph{Recall:} say the Italian for roll of film, then the Italian for to develop, then say \emph{il giardinaggio} again
\end{itemize}

\begin{tcolorbox}[breakable,colback=teal!4,colframe=teal!35!black,title={Before you move on}]
What does il giardinaggio mean? (\textquotedblleft{}Gardening\textquotedblright{}.) Say it once more.
\end{tcolorbox}
\section[seminare]{seminare — to sow}

\label{lesson:IT-C507-seminare}



\begin{quote}
Before the new one: say the Italian for to develop, then the Italian for gardening.
\end{quote}

\subsection*{You'll want to know: seminare}
\textbf{seminare} — \textquotedblleft{}to sow\textquotedblright{}.

\begin{quote}
In primavera semino i pomodori.
\end{quote}

\begin{quote}
\emph{In spring I sow tomatoes.}
\end{quote}

Built on \emph{il seme} \textquotedblleft{}seed\textquotedblright{}.

\begin{cousinweb}
\emph{seminare} is inherited from Latin \emph{sēmināre} \textquotedblleft{}to sow\textquotedblright{} (EtymDB), from \emph{sēmen} \textquotedblleft{}a seed\textquotedblright{}.
\end{cousinweb}

\begin{grammarlens}[title={What you've built}]
One more word for telling, explaining and describing at B1.
\end{grammarlens}

\subsection*{Your turn}
\begin{itemize}
\raggedright
  \item \emph{Say it:} \emph{seminare}
  \item \emph{Say it:} \emph{seminare}, once more
  \item \emph{Say it:} a sentence of your own with \emph{seminare}
  \item \emph{Recall:} say the Italian for to develop, then the Italian for gardening, then say \emph{seminare} again
\end{itemize}

\begin{tcolorbox}[breakable,colback=teal!4,colframe=teal!35!black,title={Before you move on}]
What does seminare mean? (\textquotedblleft{}To sow\textquotedblright{}.) Say it once more.
\end{tcolorbox}
\section[annaffiare]{annaffiare — to water (plants)}

\label{lesson:IT-C507-annaffiare}



\begin{quote}
Before the new one: say the Italian for gardening, then the Italian for to sow.
\end{quote}

\subsection*{You'll want to know: annaffiare}
\textbf{annaffiare} — \textquotedblleft{}to water (plants)\textquotedblright{}.

\begin{quote}
Chi annaffia le piante quando siamo in vacanza?
\end{quote}

\begin{quote}
\emph{Who waters the plants when we are on holiday?}
\end{quote}

Also spelt \emph{innaffiare}. The watering can is \emph{l'annaffiatoio}.

\begin{cousinweb}
\emph{annaffiare} is inherited from Vulgar Latin \emph{inafflāre} \textquotedblleft{}to blow onto, to sprinkle\textquotedblright{} (EtymDB), from Latin \emph{afflāre}, built on \emph{flāre} \textquotedblleft{}to blow\textquotedblright{}.
\end{cousinweb}

\begin{grammarlens}[title={What you've built}]
One more word for telling, explaining and describing at B1.
\end{grammarlens}

\subsection*{Your turn}
\begin{itemize}
\raggedright
  \item \emph{Say it:} \emph{annaffiare}
  \item \emph{Say it:} \emph{annaffiare}, once more
  \item \emph{Say it:} a sentence of your own with \emph{annaffiare}
  \item \emph{Recall:} say the Italian for gardening, then the Italian for to sow, then say \emph{annaffiare} again
\end{itemize}

\begin{tcolorbox}[breakable,colback=teal!4,colframe=teal!35!black,title={Before you move on}]
What does annaffiare mean? (\textquotedblleft{}To water (plants)\textquotedblright{}.) Say it once more.
\end{tcolorbox}
\section[il rastrello]{il rastrello — rake}

\label{lesson:IT-C507-rastrello}



\begin{quote}
Before the new one: say the Italian for to sow, then the Italian for to water.
\end{quote}

\subsection*{You'll want to know: il rastrello}
\textbf{il rastrello} — \textquotedblleft{}rake\textquotedblright{}.

\begin{quote}
Raccolgo le foglie con il rastrello.
\end{quote}

\begin{quote}
\emph{I gather the leaves with the rake.}
\end{quote}

\emph{Rastrellare} — to rake.

\begin{cousinweb}
\emph{rastrello} is inherited from Latin \emph{rāstellus} \textquotedblleft{}a small rake\textquotedblright{} (EtymDB), a little form of \emph{rāstrum}.
\end{cousinweb}

\begin{grammarlens}[title={What you've built}]
One more word for telling, explaining and describing at B1.
\end{grammarlens}

\subsection*{Your turn}
\begin{itemize}
\raggedright
  \item \emph{Say it:} \emph{il rastrello}
  \item \emph{Say it:} \emph{il rastrello}, once more
  \item \emph{Say it:} a sentence of your own with \emph{il rastrello}
  \item \emph{Recall:} say the Italian for to sow, then the Italian for to water, then say \emph{il rastrello} again
\end{itemize}

\begin{tcolorbox}[breakable,colback=teal!4,colframe=teal!35!black,title={Before you move on}]
What does il rastrello mean? (\textquotedblleft{}Rake\textquotedblright{}.) Say it once more.
\end{tcolorbox}
\section[zappare]{zappare — to hoe, to dig}

\label{lesson:IT-C507-zappare}



\begin{quote}
Before the new one: say the Italian for to water, then the Italian for rake.
\end{quote}

\subsection*{You'll want to know: zappare}
\textbf{zappare} — \textquotedblleft{}to hoe, to dig\textquotedblright{}.

\begin{quote}
Il nonno zappa l'orto ogni mattina.
\end{quote}

\begin{quote}
\emph{Grandpa hoes the vegetable garden every morning.}
\end{quote}

The tool is \emph{la zappa}.

\begin{cousinweb}
\emph{zappare} is built on \emph{zappa} \textquotedblleft{}hoe\textquotedblright{} (EtymDB), usually traced to Late Latin \emph{sappa}; the deeper origin is not settled.
\end{cousinweb}

\begin{grammarlens}[title={What you've built}]
One more word for telling, explaining and describing at B1.
\end{grammarlens}

\subsection*{Your turn}
\begin{itemize}
\raggedright
  \item \emph{Say it:} \emph{zappare}
  \item \emph{Say it:} \emph{zappare}, once more
  \item \emph{Say it:} a sentence of your own with \emph{zappare}
  \item \emph{Recall:} say the Italian for to water, then the Italian for rake, then say \emph{zappare} again
\end{itemize}

\begin{tcolorbox}[breakable,colback=teal!4,colframe=teal!35!black,title={Before you move on}]
What does zappare mean? (\textquotedblleft{}To hoe, to dig\textquotedblright{}.) Say it once more.
\end{tcolorbox}
